\section{Arquitectura del sistema}

Voctomix 2.0 se articula en torno a dos procesos principales que se comunican por red,
heredados de la versión original, a los que se añade un conjunto de servicios de soporte
que amplían sus capacidades: telemetría, gestión de audio y contenerización.

\subsection{voctocore --- Núcleo multimedia}

\texttt{voctocore} es el servidor central del sistema: gestiona todo el procesamiento
multimedia, acepta las fuentes de vídeo entrantes, realiza la composición y expone los
resultados en distintos puertos TCP. Se inicializa mediante \texttt{voctocore/voctocore.py},
que instancia en orden el \texttt{Pipeline} central (grafo GStreamer), el
\texttt{ControlServer} TCP y el proveedor de tiempo de red GStreamer
(\texttt{NetTimeProvider}, puerto UDP 9998). El \texttt{Pipeline} construye
dinámicamente la cadena GStreamer a partir de la configuración INI: crea una instancia
de \texttt{AVSource} por cada fuente declarada, instancia el \texttt{VideoMix} y el
\texttt{AudioMix}, conecta el \texttt{StreamBlanker} y abre los puertos de salida.

\subsection{voctogui --- Interfaz gráfica de control}

\texttt{voctogui} es el cliente de control: implementado en PyGTK3, proporciona al
realizador la interfaz visual para operar el sistema en tiempo real. Se conecta al
núcleo mediante un socket TCP no bloqueante gestionado con \texttt{GObject.io\_add\_watch},
de forma que los eventos de red no bloquean el hilo principal de GTK. La GUI recibe
los flujos de previsualización de cada fuente en formato JPEG comprimido (puertos
14000--14005) para minimizar el consumo de ancho de banda, mientras que el monitor
principal recibe la mezcla final en bruto (puerto 11000).

\subsection{Diagrama de puertos del sistema}

La figura~\ref{fig:diagrama_sistema} muestra la arquitectura completa del sistema,
con los bloques principales y los flujos de datos entre ellos: las cámaras y fuentes
de señal auxiliares alimentan el núcleo voctocore, que genera la mezcla. El resultado
pasa por el stream blanker antes de salir por el puerto de streaming. La GUI se conecta
al núcleo para control y visualización. RabbitMQ y el servicio de telemetría completan
el sistema recopilando los eventos de estado.

% --- Figura principal ---
% \begin{figure}[H]
%   \centering
%   \includegraphics[width=0.95\textwidth]{figuras/diagrama_sistema_completo.png}
%   \caption{Arquitectura de Voctomix 2.0: fuentes, núcleo, salidas, GUI, RabbitMQ
%            y servicio de telemetría, todo ello en contenedor Docker.}
%   \label{fig:diagrama_sistema}
% \end{figure}

La tabla~\ref{tab:puertos_sistema} recoge los puertos TCP/UDP utilizados por voctocore
y su función dentro del sistema.

\begin{table}[H]
  \centering
  \caption{Puertos del sistema Voctomix 2.0.}
  \label{tab:puertos_sistema}
  \begin{tabular}{|c|l|l|}
    \hline
    \textbf{Puerto} & \textbf{Protocolo} & \textbf{Función} \\
    \hline
    9998     & UDP & Sincronización de reloj de red GStreamer \\
    9999     & TCP & Protocolo de control (GUI, scripts, notario) \\
    10000    & TCP & Entrada cámara 1 (MKV/raw) \\
    10001    & TCP & Entrada cámara 2 (MKV/raw) \\
    10002    & TCP & Entrada cámara 3 (MKV/raw) \\
    10003    & TCP & Entrada cámara 4 (MKV/raw) \\
    10004    & TCP & Entrada fuente break (vídeo de pausa) \\
    10005    & TCP & Entrada fuente intro (vídeo introductorio) \\
    11000    & TCP & Salida mezcla principal en bruto \\
    12000    & TCP & Salida previsualización de mezcla (GUI) \\
    14000--14005 & TCP & Salidas JPEG comprimidas por fuente (GUI) \\
    15000    & TCP & Salida de streaming (tras stream blanker) \\
    17000    & TCP & Entrada vídeo stream blanker PAUSE \\
    17001    & TCP & Entrada vídeo stream blanker NOSTREAM \\
    18000    & TCP & Entrada audio (música de fondo) \\
    8080     & HTTP & API REST del servicio de telemetría \\
    \hline
  \end{tabular}
\end{table}

% --- Figura adicional ---
% \begin{figure}[H]
%   \centering
%   \includegraphics[width=0.95\textwidth]{figuras/diagrama_puertos.png}
%   \caption{Diagrama de puertos de Voctomix 2.0: numeración de cada conexión entre
%            fuentes, núcleo y consumidores.}
%   \label{fig:diagrama_puertos}
% \end{figure}

\subsection{Protocolo de control}

La comunicación entre la GUI, el servicio de telemetría y el núcleo se realiza
mediante un protocolo de texto plano sobre TCP (puerto 9999). Cada comando tiene la
forma \texttt{<verbo>\_<objeto> [parámetros]}. Por ejemplo:

\begin{itemize}
  \item \texttt{set\_video\_a cam2} --- selecciona cam2 como fuente A de vídeo.
  \item \texttt{set\_composite\_mode pip} --- activa el composite Picture in Picture.
  \item \texttt{set\_stream\_blank pause} --- activa el estado PAUSE en el stream.
\end{itemize}

El servidor transmite por \textit{broadcast} cada cambio de estado a todos los clientes
conectados, de forma que múltiples instancias de voctogui o scripts de automatización
mantienen siempre una vista coherente del sistema.
